\clearpage
\VerdierPhysicalPageAnchor{0199}%
\begingroup\small\setlength{\parskip}{0pt}\setlength{\emergencystretch}{1.5em}
\setlength{\abovedisplayskip}{4pt plus 2pt minus 2pt}%
\setlength{\belowdisplayskip}{4pt plus 2pt minus 2pt}%

\VerdierSemanticAnchor{ch3d-2-3-2-exact-sequences}
and exact sequences
\[
0\longrightarrow H^{-q}(X_{q-1})\longrightarrow
\frac{Z^{-q}}{\operatorname{Im}(Y^{-q-1}\longrightarrow Z^{-q})}
\longrightarrow
\frac{Z^{-q}}{\operatorname{Ker}(Z^{-q}\longrightarrow L^{-q+1})}
\longrightarrow0,
\]
\[
0\longrightarrow
\frac{Z^{-q}}{\operatorname{Ker}(Z^{-q}\longrightarrow L^{-q+1})}
\longrightarrow L^{-q+1}\longrightarrow
\frac{L^{-q+1}}{\operatorname{Im}(Z^{-q}\longrightarrow L^{-q+1})}
\longrightarrow0,
\]
\[
0\longrightarrow H^{-q+1}(X_{q-1})\longrightarrow
\frac{L^{-q+1}}{\operatorname{Im}(Z^{-q}\longrightarrow L^{-q+1})}
\longrightarrow
\frac{L^{-q+1}}{\operatorname{Ker}(L^{-q+1}\longrightarrow L^{-q+2})}
\longrightarrow0.
\]

\VerdierSemanticAnchor{ch3d-2-3-2-mq-object}
Since the complex $X_{q-1}$ is quasi-isomorphic to $Y$, the objects
$H^i(X_{q-1})$ belong to $M$. It follows from the exact sequences above,
the induction hypothesis 2), and the exactness properties of $M^\wedge$
(2.2.4) that the object
\[
Z^{-q}/\operatorname{Im}(Y^{-q-1}\longrightarrow Z^{-q})
\]
belongs to $M^\wedge$. Consequently, there is a morphism

\VerdierSemanticAnchor{ch3d-2-3-2-v}
\[
v\colon L^{-q}\longrightarrow Z^{-q}
\]
such that $L^{-q}$ belongs to $M$ and the composite morphism

\VerdierSemanticAnchor{ch3d-2-3-2-composite}
\[
L^{-q}\xrightarrow{v}Z^{-q}\longrightarrow
Z^{-q}/\operatorname{Im}(Y^{-q-1}\longrightarrow Z^{-q})
\]
is an epimorphism. Set

\VerdierSemanticAnchor{ch3d-2-3-2-xq}
\[
X_q=X_{q-1}(v,-q),
\qquad
\dot u_q=p(v,-q)\quad(2.1.1).
\]

\VerdierSemanticAnchor{ch3d-2-3-2-verify}
It follows from (2.1.2) that $u_q$ is a quasi-isomorphism, which proves
property 3). It remains to verify property 2), that is, to verify that the
kernel and image of $L^{-q}\longrightarrow L^{-q+1}$ belong to $M^\wedge$;
this follows immediately from the exact sequences

\VerdierSemanticAnchor{ch3d-2-3-2-exact-sequences-2}
\[
0\longrightarrow\operatorname{Im}(L^{-q}\longrightarrow L^{-q+1})
\longrightarrow\operatorname{Ker}(L^{-q+1}\longrightarrow L^{-q+2})
\longrightarrow H^{-q+1}(X_q)\longrightarrow0,
\]
\[
0\longrightarrow\operatorname{Ker}(L^{-q}\longrightarrow L^{-q+1})
\longrightarrow L^{-q}\longrightarrow
\operatorname{Im}(L^{-q}\longrightarrow L^{-q+1})\longrightarrow0.
\]

\VerdierSemanticAnchor{ch3d-2-3-2-conclusion}
This completes the proof of the proposition.

\VerdierSemanticAnchor{ch3d-2-3-3}
\medskip\noindent\textbf{Proposition 2.3.3.}
\textit{Let $M\subset M'$ be two sets of objects of $\mathcal A$ such that
$M$ satisfies properties $(E_4)$ and $(E_5)$, $M'$ satisfies properties
$(E_1)$ and $(E_4)$, and every object of $M'$ is a quotient of an object of
$M$. Let $Y\in\operatorname{Ob}\operatorname{comp}(\mathcal A)$ be a complex
such that:}

\noindent\textit{1) The components $Y^i$, $i\in\mathbf Z$, of $Y$ belong to
$M'$.}
\endgroup
